% Created 2026-02-17 mar 21:12
% Intended LaTeX compiler: pdflatex
\documentclass[11pt]{article}
\usepackage[utf8]{inputenc}
\usepackage[T1]{fontenc}
\usepackage{graphicx}
\usepackage{longtable}
\usepackage{wrapfig}
\usepackage{rotating}
\usepackage[normalem]{ulem}
\usepackage{amsmath}
\usepackage{amssymb}
\usepackage{capt-of}
\usepackage{hyperref}
\usepackage{minted}

\usepackage{parskip}
\usepackage[utf8]{inputenc} %% For unicode chars
\usepackage{placeins}
\usepackage[margin=2.50cm]{geometry}
\usepackage[T1]{fontenc}
\usepackage{mathpazo}
\linespread{1.05}
\usepackage[scaled]{helvet}
\usepackage{courier}
\hypersetup{colorlinks=true,linkcolor=blue}
\RequirePackage{fancyvrb}
\DefineVerbatimEnvironment{verbatim}{Verbatim}{fontsize=\small,formatcom = {\color[rgb]{0.5,0,0}}}
\usepackage{mdframed}
\BeforeBeginEnvironment{minted}{\begin{mdframed}}
\AfterEndEnvironment{minted}{\end{mdframed}}
\setcounter{secnumdepth}{3}
\author{José Manuel Sánchez Álvarez}
\date{}
\title{Explorando redes neuronales}
\hypersetup{
 pdfauthor={José Manuel Sánchez Álvarez},
 pdftitle={Explorando redes neuronales},
 pdfkeywords={},
 pdfsubject={},
 pdfcreator={Emacs 27.1 (Org mode 9.6.18)}, 
 pdflang={Spanish}}
\begin{document}

\maketitle
\setcounter{tocdepth}{3}
\tableofcontents

\setlength\parindent{10pt}

\section{Tarea de Clase: Explorando TensorFlow Playground}
\label{sec:org68ff327}

\subsection{Objetivo de la Tarea:}
\label{sec:org305690d}
Utilizar \href{https://playground.tensorflow.org/\#activation=tanh\&batchSize=10\&dataset=circle\&regDataset=reg-plane\&learningRate=0.03\&regularizationRate=0\&noise=0\&networkShape=4,2\&seed=0.19905\&showTestData=false\&discretize=false\&percTrainData=50\&x=true\&y=true\&xTimesY=false\&xSquared=false\&ySquared=false\&cosX=false\&sinX=false\&cosY=false\&sinY=false\&collectStats=false\&problem=classification\&initZero=false\&hideText=false}{TensorFlow Playground}, un simulador de redes neuronales, para
entrenar clasificadores binarios, ajustar la arquitectura del modelo y
sus hiperparámetros, y obtener una comprensión intuitiva del
funcionamiento de las redes neuronales y el impacto de sus
hiperparámetros.

Alternativas a Tensorflow Playground:
\begin{itemize}
\item \href{https://deeperplayground.org/}{Deeper Playgound}
\item \href{https://playground.scienxlab.org/}{ScienceXlab}
\end{itemize}

\subsection{Descripción de la Tarea:}
\label{sec:org89889b4}
Esta tarea está diseñada para permitirte experimentar directamente con
redes neuronales sin necesidad de escribir código. A través de
TensorFlow Playground, explorarás cómo diferentes configuraciones
afectan el aprendizaje y la clasificación en redes neuronales.

\subsubsection{Instrucciones:}
\label{sec:orgc87c837}

\begin{enumerate}
\item \textbf{\textbf{Entrenamiento de la Red Neuronal Predeterminada:}}
\begin{itemize}
\item Accede a TensorFlow Playground y entrena la red neuronal
predeterminada haciendo clic en el botón "Ejecutar". Observa cómo
la red encuentra una solución para la tarea de clasificación.
\item Examina los patrones aprendidos por las neuronas en cada capa
oculta. Reflexiona sobre cómo las capas combinan patrones simples
para formar patrones más complejos.
\end{itemize}

\item \textbf{\textbf{Exploración de Funciones de Activación:}}
\begin{itemize}
\item Cambia la función de activación de "tanh" a "ReLU" y entrena la
red nuevamente. Nota la diferencia en el tiempo de entrenamiento
y la forma de los límites de decisión.
\end{itemize}

\item \textbf{\textbf{Análisis del Riesgo de Mínimos Locales:}}
\begin{itemize}
\item Modifica la arquitectura de la red para tener una sola capa
oculta con tres neuronas. Entrénala varias veces, restableciendo
los pesos con el botón "Reiniciar". Documenta la variabilidad en
el tiempo de entrenamiento y la ocurrencia de mínimos locales.
\end{itemize}

\item \textbf{\textbf{Redes Neuronales Demasiado Pequeñas:}}
\begin{itemize}
\item Reduce el número de neuronas a dos y evalúa la capacidad de la
red para encontrar una solución adecuada. Discute los casos de
subajuste.
\end{itemize}

\item \textbf{\textbf{Redes Neuronales Suficientemente Grandes:}}
\begin{itemize}
\item Incrementa el número de neuronas a ocho y entrena la red varias
veces. Observa la consistencia en el rendimiento y la
velocidad. Reflexiona sobre la teoría de redes neuronales y
mínimos locales.
\end{itemize}

\item \textbf{\textbf{Desafíos de Gradientes Desvanecidos en Redes Profundas:}}
\begin{itemize}
\item Selecciona el conjunto de datos "espiral" y configura la red con
cuatro capas ocultas de ocho neuronas cada una. Analiza el
impacto en el tiempo de entrenamiento y la incidencia de mesetas.
\end{itemize}

\item \textbf{\textbf{Experimentación Adicional:}}
\begin{itemize}
\item Dedica aproximadamente una hora a experimentar con diferentes
configuraciones de hiperparámetros y arquitecturas. Intenta
obtener una comprensión intuitiva de la influencia de cada
parámetro en el aprendizaje de la red.
\end{itemize}
\end{enumerate}
\end{document}
